\documentclass{article}
\usepackage{mathtools} 
\usepackage{fontspec}
\usepackage[UTF8]{ctex}
\usepackage{amsthm}
\usepackage{mdframed}
\usepackage{xcolor}
\usepackage{amssymb}
\usepackage{amsmath}
\usepackage{tikz}
\usepackage{pgfplots}
\usepgfplotslibrary{polar}
% \pgfplotsset{compat=1.18}


% 定义新的带灰色背景的说明环境 zremark
\newmdtheoremenv[
  backgroundcolor=gray!10,
  % 边框与背景一致，边框线会消失
  linecolor=gray!10
]{zremark}{说明}

% 通用矩阵命令: \flexmatrix{矩阵名}{元素符号}{行数}{列数}
\newcommand{\flexmatrix}[4]{
  \[
  #1 = \begin{pmatrix}
    #2_{11}     & #2_{12}     & \cdots & #2_{1#4}   \\
    #2_{21}     & #2_{22}     & \cdots & #2_{2#4}   \\
    \vdots      & \vdots      & \ddots & \vdots     \\
    #2_{#31}    & #2_{#32}    & \cdots & #2_{#3#4}
  \end{pmatrix}
  \]
}

% 简化版命令（默认矩阵名为A，元素符号为a）: \quickmatrix{行数}{列数}
\newcommand{\quickmatrix}[2]{\flexmatrix{A}{a}{#1}{#2}}

\begin{document}
\title{注释 10.3}
\author{张志聪}
\maketitle

\begin{zremark}
  定理10.1中，如何利用三角不等式得到
  \begin{align*}
    |\sigma_i| \leq |y'(\eta_i) - y'(\xi_i)|, i = 1, 2, \cdots, n
  \end{align*}
\end{zremark}

% 我们有
% \begin{align*}
%   \sigma_i = \sqrt{[x'(\xi_i)]^2 + [y'(\eta_i)]^2} - \sqrt{[x'(\xi_i)]^2 + [y'(\xi_i)]^2}
% \end{align*}
% 于是
% \begin{align*}
%   |\sigma_i| & = \left| \sqrt{[x'(\xi_i)]^2 + [y'(\eta_i)]^2} - \sqrt{[x'(\xi_i)]^2 + [y'(\xi_i)]^2} \right| \\
%              & \leq \left| \sqrt{[x'(\xi_i)]^2 + [y'(\eta_i)]^2 - [x'(\xi_i)]^2 - [y'(\xi_i)]^2} \right|
% \end{align*}
todo

\begin{zremark}
  关于$y = f(x)$的曲率公式的推导。（对书中$dx$类的表示方式不喜欢。）
\end{zremark}

任意$x_0$点处的导数为$f'(x_0)$,
在$x_0$处的角度设为$\alpha$，
由导数和斜率的关系，我们有
\begin{align*}
  \alpha = arc tan f'(x_0)
\end{align*}
$\Delta s \to 0, \Delta x \to 0$
于是曲率定义可知
\begin{align*}
  K & = \left|\lim\limits_{x \to x_0} \frac{\Delta \alpha}{\Delta s}\right|                                   \\
    & = \left|\lim\limits_{x \to x_0} \frac{\frac{\Delta \alpha}{\Delta x}}{\frac{\Delta s}{\Delta x}}\right| \\
\end{align*}
因为
\begin{align*}
  \lim\limits_{x \to x_0} \frac{\Delta \alpha}{\Delta x}
   & = \frac{1}{1 + f'^2(x_0)} f''(x_0)
\end{align*}
注：这里是对$arc tan f'(x)$在$x_0$处的导数。

由弧长公式可知
\begin{align*}
  \Delta s & = \int_{x_0}^{x} \sqrt{1 + f'^2(t)} dt
\end{align*}
于是由微积分基本定理可知
\begin{align*}
  \lim\limits_{x \to x_0} \frac{\Delta s}{\Delta x}
   & = \sqrt{1 + f'^2(x_0)}
\end{align*}

综上
\begin{align*}
  K & = \left|\frac{\frac{1}{1 + f'^2(x_0)} f''(x_0)}{\sqrt{1 + f'^2(x_0)}}\right| \\
    & = \frac{|f''(x_0)|}{(1 + f'^2(x_0))^{3/2}}
\end{align*}
由$x_0$的任意性可知，命题得证。

类似地，可以得到参数方程的式子。

\end{document}